% =====================================================================
% Section 1 -- Introduction. Written in step 6.7 (after the results).
% Inputs: claims of section 3 of the notes; notes 4.9 points 8 and 12;
% reading of Pang et al. (2015) in step 6.2.
% Tone (phase 4.3): not "we show for the first time that position
% matters" but "we quantify in closed, dimensionless form an effect that
% practice knows from isolated cases".
% =====================================================================
\section{Introduction}\label{sec:intro}

% (1) Problem and practice. Start-up and stop transients in long belts;
% the gravity take-up keeps the slack-side tension; its location is a
% design choice made from topography and "belt dynamic behaviour".
% Cite: harrison1983, cornet2002, nordell1997zisco, lodewijks2013
% (take-up next to the drive keeps the exit tension constant), cdi2002.
\TODO{Problem and practice.}

% (2) Literature. Two wave speeds (nordell1984, lodewijks2002, cdi1998);
% take-up as an embedded mass (harrison1983 Fig. 2, sakharwade2019);
% song2012 (2:1 written, mass neglected, take-up moved to the tail);
% pang2015 and li2018 (single span, end mass that shares the rigid-body
% acceleration of the belt; li2018 applies it to a tail take-up);
% gao2026 (mass and pre-tension confused); discrete models with 2:1
% kinematics (lodewijks1996, kulinowski2014, song2003).
\TODO{Literature.}

% (3) Gap: no closed-form treatment of the loop with the take-up as a
% 2:1 pulley at an arbitrary position; the literature does not use the
% free end, and the result explains three loose observations
% (Harrison's period independent of the mass; spectrum change when the
% take-up is moved to the tail; return strand at rest for one wave
% transit).
\TODO{Gap.}

% (4) "Long conveyor" in the dynamic sense: take-up tension small against
% the weight of the return strand (notes 4.9 point 12; 4.26).
% Cite suchorab2025longdistance (no strict definition in the literature).
\TODO{Long conveyor in the dynamic sense.}

% (5) Contributions and outline.
\TODO{Contributions and outline of the paper.}
